% Generated by Listas/make_statements.py from pratica-lecture-06.md, item 18. Do not edit.
\textbf{Source.} Tadelis (2013), \emph{Game Theory: An Introduction}, Exercise 10.6 ``Regulating Medications'', p.~217 {[}PDF 236{]}. Reproduced verbatim. Answer in the source: none in the book (the preface, p.~xiii, says about half are solved on Princeton UP's website; not checked).



\begin{quote}

\textbf{Regulating Medications:} Consider a firm (player 1) that produces a unique drug that is used by a consumer (player 2). This drug is regulated by the government so that its price is \(p = 6\). This price is fixed, but the quality of the drug depends on the manufacturing procedure. The ``good'' (\(G\)) manufacturing procedure costs 4 to the firm and yields a value of 9 to the consumer. The ``bad'' (\(B\)) manufacturing procedure costs 0 to the firm and yields a value of 4 to the consumer. The consumer can choose whether or not to buy at the price \(p\), and this decision must be made before the actual manufacturing procedure is revealed. However, after consumption, the true quality of the drug is revealed to the consumer. The choice of manufacturing procedure, and thus of the cost of production, is made before the firm knows whether the consumer will buy or not.



\begin{enumerate}

\def\labelenumi{\alph{enumi}.}

\item

  Draw the game tree and the matrix of this game, and find all the Nash equilibria.

\item

  Now assume that the game is repeated twice. (The consumer learns the quality of the product in each period only if he consumes.) Assume that each player tries to maximize the (nondiscounted) sum of his stage payoffs. Find all the subgame-perfect equilibria of this game.

\item

  Now assume that the game as repeated infinitely many times. Assume that each player tries to maximize the discounted sum of his or her stage payoffs, where the discount rate is \(\delta \in (0, 1)\). What is the range of discount factors for which the good manufacturing procedure will be used as part of a subgame-perfect equilibrium?

\item

  Consumer advocates are pushing for a lower price for the drug, say 5. The firm wants to approach the Food and Drug Administration and argue that if the regulated price is decreased to 5 then this may have dire consequences for both consumers and the firm. Can you make a formal argument, using the given parameters, to support the firm? What about the consumers?

\end{enumerate}

\end{quote}

